\chapter{Hyperprolactinemia}
\index{Hyperprolactinemia}
\label{sec:Hyperprolactinemia}

\section{Overview}
Hyperprolactinemia is the most common pituitary hormone excess disorder, characterized by elevated serum prolactin levels.
\section{Causes}
Causes include prolactin-secreting pituitary adenomas (prolactinomas, most common), medications (antipsychotics, metoclopramide, domperidone, verapamil), hypothalamic/pituitary stalk compression, primary hypothyroidism (increased TRH stimulates prolactin), chronic kidney failure, and chest wall trauma.     
\section{Risk Factors}

\begin{itemize}[noitemsep]
  \item Genetic predisposition and family history
  \item Advancing age
  \item Lifestyle factors (diet, exercise, smoking, alcohol)
  \item Environmental exposures
  \item Underlying medical conditions
  \item Medication use
\end{itemize}
\section{Signs and Symptoms}

\subsection{Common symptoms}
\begin{itemize}[noitemsep]
  \item Clinical features in women include galactorrhea, oligomenorrhea or amenorrhea, infertility, and decreased libido. In men: decreased libido, erectile dysfunction, infertility, gynecomastia, and galactorrhea (rare). Mass effects from macroprolactinomas (>1 cm) include headache, visual field defects (bitemporal hemianopia), and hypopituitarism. 
  \end{itemize}

\subsection{Emergency warning signs}
\begin{itemize}[noitemsep]
  \item Severe or worsening symptoms of hyperprolactinemia require immediate medical evaluation
\end{itemize}
\section{Progression and Stages}
The condition may progress through different stages of severity over time. Early stages often involve mild or subtle symptoms that may be overlooked. Moderate stages involve more noticeable symptoms affecting daily function. Advanced stages may involve significant impairment and complications.
\section{Complications}
\begin{itemize}[noitemsep]
  \item Disease progression leading to worsening symptoms
  \item Secondary infections or injuries
  \item Side effects of treatment
  \item Reduced quality of life
  \item Emotional and psychological impact
\end{itemize}
\section{Diagnosis}
Doctors diagnose this based on confirmed elevated serum prolactin (macroprolactin excluded), MRI pituitary with contrast to identify adenoma, and thyroid function tests to exclude hypothyroidism. Comprehensive medical history and physical examination Laboratory tests to assess organ function and rule out other conditions Imaging studies to visualize affected structures Specialized testing depending on the affected system Consultation with relevant specialists Monitoring of disease progression over time

\begin{itemize}[noitemsep]
  \item Comprehensive medical history and physical examination
  \item Laboratory tests to assess organ function
  \item Imaging studies to visualize affected structures
  \item Specialized testing depending on the affected system
  \item Consultation with relevant specialists
\end{itemize}
\section{Prevention}
\begin{itemize}[noitemsep]
  \item Maintain a healthy lifestyle with balanced diet and exercise
  \item Avoid known risk factors
  \item Attend regular medical check-ups for early detection
  \item Follow recommended screening guidelines
\end{itemize}
\section{Treatment Options}
Treatment options include dopamine agonists (cabergoline preferred over bromocriptine due to better efficacy and tolerability) as first-line therapy for symptomatic hyperprolactinemia and prolactinomas, and transsphenoidal surgery for dopamine agonist-resistant or apoplectic macroprolactinomas. Medications to manage symptoms and address underlying mechanisms Lifestyle modifications and self-care measures Physical therapy and rehabilitation Surgical intervention when indicated Regular monitoring and follow-up care Patient education and support services

\begin{itemize}[noitemsep]
  \item Medications to manage symptoms and address underlying mechanisms
  \item Lifestyle modifications and self-care measures
  \item Physical therapy and rehabilitation
  \item Surgical intervention when indicated
  \item Regular monitoring and follow-up care
\end{itemize}
\section{Living with It}
\begin{itemize}[noitemsep]
  \item Follow your treatment plan as prescribed
  \item Maintain regular follow-up with your healthcare provider
  \item Make necessary lifestyle adjustments
  \item Seek support from family, friends, or support groups
  \item Stay informed about your condition
\end{itemize}
\section{Outlook}
\begin{itemize}[noitemsep]
  \item Most people recover well
  \item Most patients achieve normal prolactin levels with dopamine agonist therapy.
\end{itemize}
